%
% 15-analytisch.tex -- Potenzreihenentwicklung einer holomorphen Funktion
%
% (c) 2026 Prof Dr Andreas Müller
%
\subsection{Potenzreihenentwicklung einer holomorphen Funktion
\label{buch:analytisch:potenzreihen:subsection:analytisch}}
Falls für eine komplexe Funktion eine konvergente Potenzreihenentwicklung
existiert, dann ist die Funktion sogar unendlich oft stetig differenzierbar.
Ob auch die Umkehrung gilt, wissen wir noch nicht.
Die Existenz einer konvergenten Potenzreihe und die Existenz einer
komplexen Ableitung sind also nach unserem aktuellen Wissensstand
nicht notwendigerweise dasselbe.

Wir wissen aber nach
Satz~\ref{buch:analytisch:potenzreihen:satz:integralanalytisch},
dass eine holomorphe Funktion, die durch das
Integral~\eqref{buch:analytisch:potenzreihen:eqn:fzintegral}
aus den Werten einer stetigen Funktion entlang einer Kurve
entsteht, analytisch ist.
Wir können diese Erkenntnis nun mit dem cauchyschen Integralsatz
zusammenbringen und zeigen, dass jede holomorphe Funktion auch
eine konvergente Potenzreihenentwicklung hat.

Sei jetzt also $f(z)$ eine holomorphe Funktion und $\gamma$ ein
geschlossener Weg, der ein Gebiet $U$ berandet und genau einmal
umläuft.
Satz~\ref{buch:analytisch:potenzreihen:satz:integralanalytisch}
sagt für $g(x)=f(z)$, dass das
Integral~\eqref{buch:analytisch:potenzreihen:eqn:fzintegral}
eine analytische Funktion ist.
Der cauchysche Integralsatz sagt andererseits, dass sie in $U$ mit
$f(z)$ übereinstimmt.
Damit haben wir den folgenden Satz gezeigt.

\begin{satz}[holomorphe Funktionen sind analytisch]
\label{buch:analytisch:satz:holomorph-analytisch}
Eine holomorphe Funktion $f(z)$ ist analytisch, die Taylor-Reihe
von $f$ im Punkt $a$ konvergiert.
Der Konvergenzradius ist das Supremum der Radien von Kreisen um $a$,
die vollständig im Definitionsbereich von $f$ enthalten sind.
\end{satz}

Die Integrale auf der rechten Seite von \eqref{buch:analytisch:eqn:reihe}
können ebenfalls mit dem Satz~\ref{buch:integration:satz:ableitungn}
ausgerechnet werden.
Wegen
\[
f^{(k)}(z)
=
\frac{k!}{2\pi i}
\oint_\gamma
\frac{f(x)}{(x-z)^{k+1}}
\,dx
\]
ergeben sie
\begin{align*}
f(z)
&=
\sum_{k=0}^{\infty}
\biggl(
\frac{1}{2\pi i}
\oint_\gamma
\frac{f(x)}
{(x-z_0)^{k+1}}
\,dx
\biggr)
(z-z_0)^k
=
\sum_{k=0}^\infty
\frac{f^{(k)}(z_0)}{k!}
(z-z_0)^k.
\end{align*}
Damit ist die gefundene Potenzreihe der Funktion $f$ die
Taylor-Reihe der Funktion $f$ im Punkt $z_0$.

Der Satz~\ref{buch:analytisch:satz:holomorph-analytisch}
hat gezeigt, dass eine holomorphe Funktion immer auch
analytisch ist.
Dies ist eine Eigenschaft, die nicht für reelle analytische
Funktionen gilt, wie im nachfolgenden Abschnitt gezeigt wird.
